% D4 — maquette d'une fiche d'activité : recto verso élève, puis page encadrant.
% Contenu : bir/fiches/seance-08-eleve.md et seance-08-encadrant.md.
\documentclass{BIRdoc}
\BIRdocument{Fiche d'activité nº 8}{Séance 8 --- Le spectre et les bandes}
\begin{document}

\BIRtitre{Je dresse ma carte du spectre}{Noms : \BIRligne[7cm] \hfill Date : \BIRligne[3cm]}

\textbf{Aujourd'hui, nous allons} parcourir le spectre avec un récepteur relié à
l'ordinateur, et noter ce que nous y trouvons.

\begin{BIRalerte}{Règle du jour}
Nous n'écoutons que les radios publiques et les radioamateurs. Pour tout le
reste, le son reste coupé : nous regardons, nous notons la fréquence et la
forme du signal.
\end{BIRalerte}

\section*{1. Je prends l'écran en main}

Réglez le récepteur sur une station de radio FM que vous connaissez.

\begin{itemize}[itemsep=5pt]
\item Sa fréquence : \BIRligne[3cm] MHz
\item Sur le spectre (la courbe), le signal forme : \BIRcase un pic \quad \BIRcase un creux
\item Sur la chute d'eau, sa trace est : \BIRcase un trait fin \quad \BIRcase une bande large
\item De gauche à droite, on lit : \BIRligne[6cm]
\item De haut en bas de la chute d'eau, on lit : \BIRligne[6cm]
\end{itemize}

\section*{2. Je cherche ces signaux}

\begin{tabularx}{\linewidth}{|P{48mm}|P{27mm}|P{27mm}|L|}
\hline\BIRentete
\BIRth{Signal} & \BIRth{Fréquence relevée} & \BIRth{Domaine : HF, VHF, UHF} & \BIRth{À la chute d'eau, il ressemble à…} \\ \hline
\BIRhaut une deuxième station FM & & & \\ \hline
\BIRhaut la radio numérique DAB+ & & & \\ \hline
\BIRhaut le relais radioamateur indiqué par l'encadrant & & & \\ \hline
\BIRhaut une balise radioamateur & & & \\ \hline
\end{tabularx}

\section*{3. Je calcule}

300 divisé par la fréquence en MHz donne la longueur d'onde en mètres.

\begin{tabularx}{\linewidth}{|L|P{40mm}|P{40mm}|}
\hline\BIRentete
\BIRth{Signal} & \BIRth{Fréquence en MHz} & \BIRth{Longueur d'onde en m} \\ \hline
\BIRhaut ma station FM & & \\ \hline
\BIRhaut le relais radioamateur & & \\ \hline
\end{tabularx}

\medskip
Le relais se trouve dans la bande radioamateur des \BIRligne[3cm]

\newpage
\section*{4. Je découvre}

Trouvez un signal qui n'est pas dans le tableau. Son coupé, décrivez-le.

\begin{itemize}[itemsep=5pt]
\item Fréquence : \BIRligne[3.5cm] \quad Domaine : \BIRligne[2.5cm]
\item Il est : \BIRcase permanent \quad \BIRcase par moments
\item Sa trace : \BIRligne[11cm]
\end{itemize}

\section*{5. Ma carte du spectre}

Placez vos signaux sur l'axe, chacun avec son nom.

\vspace{32mm}
{\sffamily\small
\setlength{\unitlength}{1mm}
\begin{picture}(170,14)
\put(0,8){\line(1,0){170}}
\multiput(0,5)(34,0){6}{\line(0,1){6}}
\put(0,0){\makebox(0,0)[bl]{80 MHz}}
\put(34,0){\makebox(0,0)[b]{100}}
\put(68,0){\makebox(0,0)[b]{150}}
\put(102,0){\makebox(0,0)[b]{200}}
\put(136,0){\makebox(0,0)[b]{300}}
\put(170,0){\makebox(0,0)[br]{450 MHz}}
\end{picture}}

\vspace{6mm}
\begin{BIRencadre}{À retenir}
\begin{itemize}[itemsep=6pt]
\item Le spectre se découpe en \BIRligne[3.5cm] ; les trois plus utiles sont
  \BIRligne[1.5cm], \BIRligne[1.5cm] et \BIRligne[1.5cm].
\item Une bande radioamateur porte le nom de sa \BIRligne[5cm]
\item Sur une chute d'eau, la fréquence se lit de \BIRligne[2.2cm] à
  \BIRligne[2.2cm] et le temps de \BIRligne[2.2cm] en \BIRligne[2.2cm]
\end{itemize}
\end{BIRencadre}

% ------------------------------------------------------------------ encadrant
\newpage
\BIRdocument{Fiche d'activité nº 8 --- encadrant}{Séance 8 --- Le spectre et les bandes}
\BIRtitre{Je dresse ma carte du spectre}{Côté encadrant --- manipulation de 50 min, groupes de 2 ou 3}

\begin{minipage}[t]{.36\linewidth}
\section*{Matériel par groupe}
\begin{itemize}
\item ordinateur, logiciel de réception SDR installé ;
\item clé SDR avec son antenne ;
\item fiche élève, crayon, calculatrice.
\end{itemize}
\end{minipage}\hfill
\begin{minipage}[t]{.60\linewidth}
\section*{La veille}
\begin{itemize}
\item Brancher chaque clé sur le poste où elle servira et vérifier la réception.
\item Relever depuis la salle les fréquences de deux stations FM, d'un
  multiplex DAB+, du relais local et d'une balise.
\item Régler le logiciel pour que spectre et chute d'eau soient visibles ensemble.
\item Prévoir un WebSDR ouvert sur le poste projeté, en secours.
\end{itemize}
\end{minipage}

\section*{Pendant la manipulation}

\begin{tabularx}{\linewidth}{P{26mm}P{48mm}L}
\BIRentete
\BIRth{Temps} & \BIRth{Étape de la fiche} & \BIRth{À surveiller} \\
0 à 10 min & 1. prise en main & que chaque groupe ait un signal à l'écran \\
10 à 30 min & 2. recherche, 3. calcul & les groupes qui restent sur la bande FM \\
30 à 45 min & 4. découverte & le son coupé hors radiodiffusion et bandes radioamateur \\
45 à 50 min & 5. carte & l'ordre des signaux sur l'axe \\
\bottomrule
\end{tabularx}

\section*{Pièges matériels}
\begin{itemize}
\item Rien à l'écran : gain à zéro, mauvaise entrée choisie, antenne débranchée.
\item Pics fixes qui restent sans antenne : parasites, à ignorer.
\item Relais muet : il n'émet que lorsqu'il est utilisé. Le radioamateur
  présent peut le déclencher pour la démonstration.
\item DAB+ absent : toutes les zones ne sont pas couvertes ; remplacer par une
  troisième station FM et le dire.
\end{itemize}

\begin{BIRencadre}{Savoir-faire à valider dans le livret}
\BIRcase lit une fréquence sur l'écran \quad
\BIRcase distingue spectre et chute d'eau \\[3pt]
\BIRcase range un signal dans HF, VHF ou UHF \quad
\BIRcase calcule une longueur d'onde
\end{BIRencadre}

\end{document}
